
\section{Gamma--Appell normalization and a single height-one coordinate}
\label{gat:section}

The preceding normal forms explain differential cancellations.  The
scalar Gamma normalization also controls the singular spectral
coordinates at every degree.  We reconcile the all-degree argument in
\cite{dpm:reports} with the height-one convention already fixed in
Section~\ref{dgp:subsec:heightkernel}.  Its remaining derivatives have
ordinary convergent defining sums; no independence or smallest possible
arithmetic basis is asserted.

Define coefficientwise polynomials and their fixed-coordinate moments by
\[
 \mathsf P_m(Y)=m![t^m]\Gamma(1-t)e^{tY},\qquad
 \mathsf S_m=\operatorname{FP}_x\int_0^1\mathsf P_m(\psi(x))\,dx.
\]
These symbols do not denote the separate Gaussian $S_6$ and $S_8$
ladder candidates.  Write
\[
 \alpha_j=[t^j]\Gamma(1+t)^{-1},\qquad
 L_d=[s]Z(1+s,1^{d-1})\quad(d\ge1),
\]
in the strict largest-index-first convention of
\eqref{dgp:eq:strict-zeta}.  Thus $L_1=-\gamma_1$, $L_2=\eta$,
$L_3=\theta$, and $L_4=\chi$.  A linear Laurent coefficient at a
pole is not an ordinary derivative at that pole.

Let $a_{m,j}(n)$ be the coefficient of $(x+n)^{-j}$ in
$\mathsf P_m(\psi(x))$, $1\le j\le m$, and set
\begin{align*}
 e_r(n)&=H_n(1^r),\qquad e_0(n)=1,\\
 \delta_m(n)&=a_{m,1}(n)+m!e_{m-1}(n),\\
 \mathsf J_m&=\sum_{n\ge1}\sum_{j=2}^m
       \frac{a_{m,j}(n)-a_{m,j}(n-1)}{(j-1)n^{j-1}},\\
 \mathsf D_m&=-\sum_{n\ge1}
       \bigl(\delta_m(n)-\delta_m(n-1)\bigr)\log n.
\end{align*}
The next result proves convergence and gives a finite description of
both corrections.

\begin{theorem}[One height-one coordinate at every degree]
\label{gat:main}
For every integer $m\ge2$,
\begin{equation}\label{gat:spectral}
 \boxed{\mathsf S_m=m!(\alpha_m-L_{m-1})
                          +\mathsf J_m+\mathsf D_m.}
\end{equation}
The higher-pole sum $\mathsf J_m$ is a finite linear combination of
admissible multiple zeta values with coefficients in the ordinary zeta
algebra.  The correction $\mathsf D_2=\mathsf D_3=0$; for $m\ge4$,
$\mathsf D_m$ is a finite linear combination of first coordinate
derivatives of admissible multiple zeta functions, of depth at most
$m-3$.  Each retained derivative has an ordinary absolutely convergent
series and a holomorphic parameter neighborhood.  The construction
retains no lower height-one Laurent coefficient.
\end{theorem}

\subsection{The tail ideal and the finite marked-word compiler}

Put $h=H_n$ and $A_r(n)=\zeta(r)-H_n^{(r)}$ for $r\ge2$.
These are strict upper tails.  Only finitely many enter a fixed pole
coefficient.

\begin{lemma}[The residue defect is in the upper-tail ideal]
\label{gat:tailideal}
The defect $\delta_m(n)$ is a polynomial in $h,A_2,A_3,\ldots$
with coefficients in the ordinary zeta algebra, and vanishes when
all the $A_r$ are set to zero.  For $m\ge4$, its degree in these
harmonic variables, each counted with degree one, is at most $m-3$.
The defects at degrees two and three are identically zero.
\end{lemma}

\begin{proof}
Use the local series \eqref{dgp:eq:pole-series}.  Multiplication by
$\Gamma(1-t)$ cancels the $-\gamma$ in its regular constant, since
\[
 \Gamma(1-t)=\exp\left(\gamma t+\sum_{r\ge2}\zeta(r)t^r/r\right).
\]
Finite residue extraction is therefore polynomial in harmonic numbers
and ordinary zeta values.  Substitute $H_n^{(r)}=\zeta(r)-A_r$.
When every upper tail is zero the local series is
$h-\gamma-\pi\cot(\pi v)$.  The closed residue polynomials in
Lemma~\ref{dgp:powers:lem:closed-residue} give, coefficientwise in $t$,
\[
 \operatorname*{Res}_{v=0}e^{-\pi t\cot(\pi v)}
                   =-\frac{\sin\pi t}{\pi}.
\]
Euler reflection turns its normalized residue generator into
\[
 -\Gamma(1-t)e^{(h-\gamma)t}\frac{\sin\pi t}{\pi}
             =-t\frac{e^{(h-\gamma)t}}{\Gamma(1+t)}.
\]
The elementary harmonic generator, specialized in exactly the same way,
is
\[
 \sum_{r\ge0}e_r(n)t^r
  =\exp\left(ht+\sum_{r\ge2}(-1)^{r-1}H_n^{(r)}t^r/r\right)
  \longmapsto\frac{e^{(h-\gamma)t}}{\Gamma(1+t)}.
\]
Thus the two residue contributions cancel.  Known relations among
even zeta values are retained here; those values are not treated as
independent indeterminates.

To obtain the sharper harmonic degree, first retain the variables
$h_r=H_n^{(r)}$.  In the simple-pole coefficient the only terms of
harmonic degrees $m-1$ and $m-2$ are, respectively,
\[
 -m h^{m-1},\qquad
 \frac{m!}{2(m-3)!}h_2 h^{m-3}\quad(m\ge3).
\]
The first uses one pole factor.  The second uses two pole factors and
one $h_2v$ factor; higher positive powers, additional pole factors,
or ordinary zeta coefficients lower the harmonic degree further.
The identical two terms occur in $-m!e_{m-1}$, so they cancel in
$\delta_m$.  Substitution of $h_r=\zeta(r)-A_r$ cannot increase
degree.  Direct extraction at $m=2,3$ gives zero.
\end{proof}

For the consecutive difference use $h=H_{n-1}$, $v=1/n$, and
$A_r=A_r(n)$.  The exact finite polynomial rule is
\begin{equation}\label{gat:difference}
 \delta_m(n)-\delta_m(n-1)
  =\delta_m(h+v,A_2,A_3,\ldots)
       -\delta_m(h,A_2+v^2,A_3+v^3,\ldots).
\end{equation}
Every surviving monomial has a factor $v^q$ with $q\ge1$.
If no upper tail remains, then $q\ge2$: a term with $q=1$ can
arise only by changing $h$, and the tail-ideal property leaves an
upper tail in that case.  Since $A_r(n)=O(n^{1-r})$, each summand
is $O(n^{-2}(1+\log n)^B)$ for some fixed $B$.
The Dirichlet series of this difference and all its parameter
derivatives consequently converge normally for $\operatorname{Re}s>-1$.

Here is the finite conversion.  Apply strict stuffle multiplication
separately to the upper-tail factors and the powers of $H_{n-1}$.
For an upper word $\mathbf a=(a_1,\ldots,a_r)$, whose entries are
at least two, put
\[
 T_{\mathbf a}(n)=\sum_{k_1>\cdots>k_r>n}
                          k_1^{-a_1}\cdots k_r^{-a_r}.
\]
The empty upper and lower words both have value one.  Their ordered
indices give the exact identity
\begin{equation}\label{gat:concatenation}
 \sum_{n\ge1}\frac{T_{\mathbf a}(n)H_{n-1}(\mathbf b)}{n^{q+s}}
                       =Z(\mathbf a,q+s,\mathbf b).
\end{equation}
If the upper word is empty, $q\ge2$; otherwise the first entry is
at least two.  Every such term is holomorphic for
$\operatorname{Re}s>-1$ by the same tail bounds.  Thus
\begin{equation}\label{gat:mark}
 -\sum_{n\ge1}\frac{T_{\mathbf a}(n)H_{n-1}(\mathbf b)\log n}{n^q}
               =\partial_{r+1}Z(\mathbf a,q,\mathbf b).
\end{equation}
The derivative marks the distinguished $n$ index.  It need not be
the first coordinate, and is not obtained by differentiating an
identity known only at integer indices.  A nonzero finite difference
replaces at least one harmonic factor by a power of $v$.  At most
$m-4$ harmonic factors remain; adjoining the distinguished index
gives depth at most $m-3$.

\begin{proof}[Proof of Theorem~\ref{gat:main}]
The scalar linear combination of Theorem~\ref{dgp:thm:all-power}
applies to $\mathsf P_m$.  Its large-circle constant is
\[
 m![t^m]\Gamma(1-t)\frac1{2\pi}\int_{-\pi}^{\pi}e^{it\theta}\,d\theta
   =m![t^m]\Gamma(1-t)\frac{\sin\pi t}{\pi t}=m!\alpha_m.
\]
Also $e_{m-1}(n)-e_{m-1}(n-1)=H_{n-1}(1^{m-2})/n$.
Hence its residue-difference series splits as
\[
 \sum_{n\ge1}\frac{a_{m,1}(n)-a_{m,1}(n-1)}{n^s}
  =-m!Z(1+s,1^{m-2})
       +\sum_{n\ge1}\frac{\delta_m(n)-\delta_m(n-1)}{n^s}.
\]
The second term is holomorphic at zero by \eqref{gat:difference},
and its linear coefficient is exactly $\mathsf D_m$.
The first supplies $-m!L_{m-1}$.  Higher-pole differences acquire
at least one power of $1/n$, so their outer exponent is at least two;
the finite harmonic-word compiler gives $\mathsf J_m$ as admissible
zeta values.  Equations \eqref{gat:concatenation}--\eqref{gat:mark}
prove every assertion about the mixed correction and its depth.
\end{proof}

\subsection{A completely ordinary Cauchy-series form}

Use the positive rational coefficients already defined in
Proposition~\ref{dgp:powers:prop:Cauchy}:
\[
 \rho_n=\frac1{(n-1)!}\int_0^1(r)_n\,dr,\qquad 0<\rho_n\le n.
\]
Its formula $L_{m-1}=\alpha_m+\sum_{n\ge1}\rho_n/n^{m+1}$
cancels the reciprocal-Gamma coefficient in \eqref{gat:spectral}.
Therefore, at every degree $m\ge2$,
\begin{equation}\label{gat:ordinary}
 \boxed{\mathsf S_m=\mathsf J_m+\mathsf D_m
                       -m!\sum_{n\ge1}\frac{\rho_n}{n^{m+1}}.}
\end{equation}
All three sums are ordinary and absolutely convergent.  The mixed
correction still contains its specified convergent derivatives;
this representation does not assert their further arithmetic elimination.

The lower cases make the normalization concrete:
\begin{align*}
 \mathsf S_2&=2(\alpha_2-L_1),\\
 \mathsf S_3&=6(\alpha_3-L_2)+3\zeta(2),\\
 \mathsf S_4&=24(\alpha_4-L_3)+12\zeta(3)+4\zeta'(3),\\
 \mathsf S_5&=120(\alpha_5-L_4)+\frac{260}{3}\zeta(4)+\mathsf D_5,\\
 \mathsf D_5&=20\zeta(2)\zeta'(2)-10\partial_2Z(2,2)
       +20(\partial_1-\partial_2)Z(3,1)-30\zeta'(4).
\end{align*}
Indeed $\mathsf J_2=0$, $\mathsf J_3=3\zeta(2)$,
$\mathsf J_4=12Z(2,1)=12\zeta(3)$, and
\[
 \mathsf J_5=60Z(2,1,1)+10Z(2,2)+10\zeta(2)^2-\frac{35}{6}\zeta(4)
             =\frac{260}{3}\zeta(4).
\]
The last equality uses the weight-four identities already used in
Section~\ref{dgp:sec:poles}.  Direct pole extraction gives
\[
 \delta_4=-4A_3,\qquad
 \delta_5=5A_2^2-20\zeta(2)A_2-20hA_3+25A_4.
\]
Their differences in \eqref{gat:difference} are
$4v^3$ and
$-10A_2v^2-20A_3v+20hv^3-30v^4+20\zeta(2)v^2$.
The marked-word rule and the free-order stuffle identity
$\partial_1Z(2,2)+\partial_2Z(2,2)=\zeta(2)\zeta'(2)-\zeta'(4)$
prove the displayed $\mathsf D_5$.  These equations agree with the
canonical raw fifth moment after its finite Stieltjes coordinate conversion.

\subsection{Transport through the differential normal form and recovery}

Let $A_m=N_m+\mathcal DT_m$ be Section~\ref{ibp:section}'s
inverse-Gamma decomposition, and put $K_m=N_m|_{u_0=0}$.
Since $\partial_{u_0}$ commutes with $\mathcal D$ and preserves the
normal space, uniqueness and the Gamma recurrence give
$\partial_{u_0}N_m=mN_{m-1}$.  Thus
$N_m=\sum_j\binom mj u_0^jK_{m-j}$.
Coefficientwise Gamma convolution yields the all-degree transport
\[
 \mathcal N\!\left(m![t^m]\frac{\Gamma(1-t)\Gamma(x)}{\Gamma(x-t)}\right)
                  =\sum_{j=0}^m\binom mj\mathsf P_j(\psi(x))K_{m-j}.
\]
Here $\mathcal N$ is extended linearly over constant Gamma coefficients.
The associated primitive is the same finite convolution of the $T_j$;
its endpoint functional is retained.  This statement supplies a polynomial
normal form, not an arithmetic evaluation of every resulting integral.

Finally $Q_0=1$, $\mathsf S_0=1$ and $\mathsf S_1=\gamma$.
The raw moments are recovered at every fixed degree by
\begin{equation}\label{gat:inverse}
 Q_m=m!\sum_{k=0}^m\frac{(-1)^{m-k}\alpha_{m-k}}{k!}\mathsf S_k.
\end{equation}
This follows by multiplying the coefficientwise normalization by
$1/\Gamma(1-t)$.  It interchanges no infinite series with a divergent
integral.  The finite input coefficients satisfy
\[
 \alpha_0=1,\qquad
 j\alpha_j=\gamma\alpha_{j-1}
          +\sum_{r=2}^j(-1)^{r+1}\zeta(r)\alpha_{j-r}.
\]
Independent finite residue and tail calculations through degree nine
check the compiler and its convergence rules.  Exact lower-degree
conversions and separate native quadrature check normalization; the
preceding proofs establish the all-degree identities.  Detailed sixth-
and seventh-degree arithmetic reductions remain separate audit targets.
